\documentclass{article}
\usepackage{amsmath}
\usepackage{amssymb}
\usepackage{graphicx} % Required for inserting images
\usepackage{booktabs}
\usepackage{algorithm}
\usepackage{algpseudocode}

\title{Degeneration Probe}
\author{Marco De Negri, Luca Sartori, Lorenzo Baggi}
\date{July 2026}

\begin{document}

\maketitle

\section{Dataset}
\label{sec:dataset}

Every dataset build in this project is a set of model completions (\emph{rollouts})
obtained by sampling a fixed pool of prompts and generating, for each prompt,
$K=10$ independent completions at temperature $0.7$, top-$p$ $0.9$, with a
hard generation cap of $4096$ tokens. All builds share the exact same prompt
sample, the same sources, the same per-source sample size, and the same
random seed ($42$), so that any difference in degeneration behaviour
observed between builds can be attributed to the model rather than to
prompt selection.

\subsection{Models}

We generate rollouts from three model checkpoints:

\begin{itemize}
    \item \textbf{Apertus-8B-Instruct} (\texttt{swiss-ai/Apertus-8B-Instruct-2509}), the publicly released instruction-tuned checkpoint.
    \item \textbf{Apertus 1.5 - capfilter-linear} (checkpoint \texttt{it8816}), an SFT checkpoint of Apertus 1.5.
    \item \textbf{Apertus 1.5 - SFT-256k} (checkpoint \texttt{4200}), a second, independently trained SFT checkpoint of Apertus 1.5.
\end{itemize}

\subsection{Domains}

Prompts are drawn from seven sources spanning mathematics, competitive
programming, general instruction-following, and medical reasoning. Five
sources are used in-domain, split into training, validation, and an
in-domain test set; the remaining two are held out entirely from training
and used only for zero-shot evaluation on unseen domains:

\begin{itemize}
    \item \emph{In-domain}: DeepMath-103K, NuminaMath-1.5, IF\_sft\_data\_verified, Llama-Nemotron (code split), AIME 2025
    \item \emph{Held-out (zero-shot)}: medical-o1-verifiable-problem, Codeforces (restricted to problems rated $\geq 2000$)
\end{itemize}

\section{Labeling}

Starting from the definition of degeneration alone, two natural detection
strategies emerge, depending on which side of the model one looks at. Looking
\emph{inside} the model, degeneration is a collapse of the next-token
probability distribution onto a single token or a small set of tokens; once
this happens, the model can become stuck cycling among that narrow set,
unable to escape back into a higher-entropy, non-degenerate continuation,
which suggests tracking the entropy of that distribution. Looking only at
\emph{what is generated}, degenerate text is repetitive, which can be
detected either by measuring how repetitive the token stream is locally
(TTR, Section~\ref{sec:ttr-labeling}) or by searching for the exact repeated
pattern (LRS, Section~\ref{sec:lrs-labeling}).

By manually inspecting a number of generations and their labels, we found
that none of these metrics is perfect on its own: each produces false
positives, and none reliably pinpoints the exact onset of degeneration. We
therefore introduced an LLM judge (Section~\ref{sec:llm-judge-labeling}) to
validate the structural metrics against something closer to ground truth.
This analysis confirmed that no single metric can be trusted blindly, so we
use the LLM judge as the final arbiter, followed by manual inspection of its
verdicts, which stays affordable since only a small fraction of rollouts are
ever flagged as degenerating.

\subsection{Entropy}
\label{sec:entropy-labeling}

Entropy measures, at each generated token position, how concentrated the
model's next-token probability distribution is: low entropy means the model
is committing to a narrow set of continuations, which is the mechanism one
would expect to underlie a degenerate loop. This signal correlates with
degeneration, but on its own it is noisy and produces many false positives.

\subsubsection{Definition}

Let $\mathbf{z}_t \in \mathbb{R}^{|\mathcal{V}|}$ denote the raw logits
produced by the model at generation step $t$, where $\mathcal{V}$ is the
vocabulary. The probability assigned to token $v \in \mathcal{V}$ is

\begin{equation}
p_t(v)
=
\frac{\exp(z_{t,v})}
{\sum_{u \in \mathcal{V}} \exp(z_{t,u})}.
\end{equation}

The per-token entropy is the Shannon entropy of this distribution:

\begin{equation}
H_t
=
-\sum_{v \in \mathcal{V}} p_t(v)\log p_t(v).
\label{eq:token-entropy}
\end{equation}

Natural logarithms are used, so entropy is measured in \emph{nats}. Entropy
in bits can be obtained as $H_t/\log 2$. Notice that $H_t$ characterizes the
entire next-token distribution at step $t$; it is not the surprisal
$-\log p_t(y_t)$ of the particular token $y_t$ that was sampled.

Crucially, Equation~\ref{eq:token-entropy} is evaluated on the model's raw
logits, over the full vocabulary, at temperature one. Temperature scaling
and top-$p$ truncation are applied only afterwards to sample the next token.
Computing entropy after these transformations would instead measure the
uncertainty of the truncated sampling distribution and would systematically
underestimate the model's original uncertainty.

For example, suppose that at one generation step the model assigns
probabilities $0.9$ and $0.1$ to the only two possible continuations. The
entropy recorded for that position is

\begin{align}
H_t
&=
-\left(0.9\log 0.9 + 0.1\log 0.1\right) \\
&\approx 0.3251 \text{ nats}.
\end{align}

By comparison, a uniform distribution over the same two tokens has entropy
$\log 2 \approx 0.6931$ nats, whereas a deterministic distribution has
entropy zero. Thus, lower entropy indicates a more predictable
continuation.

\subsubsection{Algorithm}

During rollout generation, the model is evaluated once at each step: the raw
logits from that forward pass are used both to compute entropy and,
separately, filtered by temperature and top-$p$ to sample the next token, so
the recorded entropy always reflects the unfiltered distribution the model
actually produced.

\begin{algorithm}[h]
\caption{Autoregressive decoding with entropy recording}
\label{alg:entropy-generation}
\begin{algorithmic}[1]
\Function{TokenEntropy}{$\mathbf{z}$}
    \State $p \gets \Call{Softmax}{\mathbf{z}}$
    \State \Return $-\sum_{v} p_v \log p_v$
\EndFunction
\Statex
\For{$t = 1, \ldots, \text{maxNewTokens}$}
    \State $\mathbf{z}_t \gets$ raw logits at step $t$
    \State $H_t \gets \Call{TokenEntropy}{\mathbf{z}_t}$ \Comment{pre-sampling}
    \State $\mathbf{z}_t' \gets \Call{ApplyTemperatureAndTopP}{\mathbf{z}_t}$
    \State $y_t \gets \Call{Sample}{\mathbf{z}_t'}$
    \State append $y_t$, $H_t$ to their respective sequences
\EndFor
\end{algorithmic}
\end{algorithm}

Consequently, each rollout stores an entropy sequence aligned one-to-one
with its generated-token sequence:

\begin{equation}
(y_1,\ldots,y_T)
\quad\longleftrightarrow\quad
(H_1,\ldots,H_T).
\end{equation}

\subsubsection{Interpretation and limitations}

Entropy is a token-level, content-agnostic signal: a drop in entropy means
the model has become confident about the next token, which can happen for
entirely legitimate reasons -- completing a well-known formula, continuing
a structured list, or closing a syntactic pattern -- and not only from a
degenerate loop. It also says nothing about \emph{which} span of text is
becoming predictable or where that span begins, unlike TTR and LRS, which
operate directly on the generated tokens. For these reasons Entropy is the
weakest and noisiest of the three metrics, confirmed empirically in
\S\ref{sec:empirical-performance}, and is used only as a complementary,
exploratory signal rather than as a labeling decision by itself.

\subsection{Type Token Ratio and Repetition Labeling}
\label{sec:ttr-labeling}

\subsubsection{Definition and computation}

The Type Token Ratio (TTR) is the fraction of distinct elements (types)
among all elements (tokens) in a sequence: $(A,A,A,A)$ has $\mathrm{TTR}=1/4$,
while $(A,B,C,D)$ has $\mathrm{TTR}=4/4=1$. Since our goal is to assign high
scores to repetitive text, we use
$\mathrm{repetition\_score}=1-\mathrm{TTR}$, so scores near $0$ indicate
diverse text and scores near $1$ indicate strong repetition.

We compute this score locally rather than over a whole completion, using a
sliding window of size $W$ over $n$-grams; in our experiments $W=256$ and
$n=2$ (bigrams). Every token position $t$ defines its own forward-looking
window $(x_t,\ldots,x_{t+W-1})$, so the result is a per-token sequence
$(r_0,\ldots,r_{L-1})$ -- one score per position, not one value per
non-overlapping chunk. Each window contains $W-n+1=255$ bigram occurrences,
and
\[
r_t = 1-\frac{\#\text{ distinct bigrams in the window}}{255}.
\]
If every bigram in the window is different, $r_t=0$; if all bigrams are
identical, $r_t=1-\tfrac1{255}\approx0.996$. Larger values of $r_t$
therefore correspond to more repetitive regions of the generated text.

\subsubsection{Algorithm}

The score is recomputed independently for each window position:

\begin{algorithm}[h]
\caption{Sliding-window repetition score}
\label{alg:ttr-sliding-window}
\begin{algorithmic}[1]
\Function{SlidingWindowRepetition}{tokens, $W$, $n$}
    \State $L \gets$ length of tokens
    \For{$t = 0, \ldots, L-1$}
        \If{$t + W > L$} \Comment{Incomplete window}
            \State $r_t \gets$ undefined
        \Else
            \State window $\gets$ tokens$[t : t+W]$
            \State $r_t \gets 1 - \dfrac{|\{\text{distinct } n\text{-grams in window}\}|}{W-n+1}$
        \EndIf
    \EndFor
    \State \Return $(r_0, \ldots, r_{L-1})$
\EndFunction
\end{algorithmic}
\end{algorithm}

\subsubsection{Undefined labels}

A label requires the full window, i.e.\ $t+W\leq L$; otherwise the
implementation assigns \texttt{NaN}. With $W=256$, the last $255$ positions
of a completion have undefined labels, and completions shorter than $256$
tokens have no valid window at all. Undefined labels are ignored during
aggregation and masked out during training.

\subsubsection{Rollout-level degeneration decision}

Each completion produces a sequence of token-level repetition labels
\[
\mathbf{r}=(r_0,\ldots,r_{L-1}).
\]
Ignoring undefined values, we summarize a rollout by its maximum valid
label,
\[
r_{\max} = \max_{t:\,r_t\text{ is defined}} r_t,
\]
and classify it as degenerating if $r_{\max}>\tau$ for a threshold $\tau$
tuned per dataset build rather than fixed a priori
(\S\ref{sec:empirical-performance}).

\subsubsection{Interpretation and limitations}

The TTR-based repetition score is a lexical, tokenizer-dependent heuristic.
It measures local \(n\)-gram diversity, not semantic degeneration directly.
Consequently, a high score can also occur in legitimate text containing
repeated mathematical expressions, structured enumerations, code templates,
or recurring formatting.

The fixed window size determines the temporal scale of the detected
repetition. Short repetitive patterns are strongly reflected in the score,
whereas repetitions separated by more than one window may not be captured.
Moreover, different tokenizers can produce different bigram sequences for the
same text.

For these reasons, the TTR-based score should be interpreted as a cheap,
continuous proxy for repetition rather than as perfect ground-truth
annotation. Its completion-level detection performance against the LLM
judge, how often it fires on naturally-terminated versus token-capped
rollouts, and how well its own onset agrees with the judge's, are reported
jointly with Entropy and LRS in \S\ref{sec:empirical-performance}.

\subsection{Longest Repeated Substring (LRS)}
\label{sec:lrs-labeling}

\subsubsection{Definition}

The repetition score of Section~\ref{sec:ttr-labeling} is a continuous,
per-token signal: it tells us how repetitive the local context is, but not
\emph{which} span of text is being repeated, nor where that span begins.
LRS is designed to answer that more specific question directly, by scanning
a whole completion for the longest span of tokens that occurs twice without
overlapping. If a model is stuck in a loop, the text it keeps repeating is,
by definition, a long substring occurring more than once; if it is not
looping, the longest substring that happens to repeat by chance is
typically short.

Formally, given a token sequence $\mathbf{x}=(x_0,\ldots,x_{L-1})$, LRS is
the largest $\ell$ for which there exist two positions $i<j$ with
$j-i\geq\ell$ and
\[
(x_i,\ldots,x_{i+\ell-1}) = (x_j,\ldots,x_{j+\ell-1}).
\]

It may seem restrictive to search for a span that repeats only \emph{twice},
rather than explicitly modeling "the same short unit repeated many times."
It is not: a short unit repeating many times in a row already satisfies this
definition, just with a longer $\ell$. If a unit $P$ repeats six times back
to back, $PPPPPP$, then the first three copies and the last three copies are
themselves identical, non-overlapping strings, a valid match of length
$3|P|$. Since LRS always keeps the longest match it finds, it reports this
combined span rather than a single copy of $P$, with no special-casing
needed for repeat counts. A second step then decomposes the matched span
into its own smallest repeating unit and counts how many times that unit
recurs across the sequence, so the result is reported as, e.g., "unit $P$
repeated 6 times" rather than as a single long match.

A degenerate loop, however, is rarely a perfectly identical chunk copied
twice: something small very often changes between repetitions, most
commonly a number. A model enumerating cases might produce
\begin{quote}
\ttfamily\raggedright N = 41: checking divisibility... N = 42: checking divisibility... N = 43: ...
\end{quote}
LRS compares token sequences position by position, so it treats the two
occurrences of the loop body as equal only where their tokens actually
match. Here the two occurrences agree everywhere except at the single token
encoding the number (\texttt{41} vs.\ \texttt{42}), so LRS cannot report the
two full occurrences as one repeat: it can only recover whichever
contiguous stretch of tokens is identical on both sides. If the number
sits at the end of the template, this means the reported repeat is
truncated right before it, understating how much of the rollout is
actually degenerate. If the number sits in the middle of the template,
as in this example, the mismatch splits the template in two and breaks the
match entirely, since no single contiguous span still matches across the
mismatch -- exactly the failure mode digit normalization (next) exists to
fix.

\subsubsection{Digit normalization}

We mitigate this by replacing each number -- regardless of how many digits
it has -- with a single generic placeholder token before matching, so that
\texttt{"N = 41"} and \texttt{"N = 42"} (or, just as well, \texttt{"N = 41"}
and \texttt{"N = 4100"}) look identical to the matcher.

\subsubsection{Algorithm}

The core routine finds the longest length $\ell$ for which a non-overlapping
repeat exists, using binary search over $\ell$: the existence of a repeat
of length $\ell$ is monotonic, since any valid repeat of length $\ell$
contains a valid repeat of every shorter length at the same two positions,
so we can binary search on $\ell$ instead of checking every possible
length. Each probe hashes every window of the candidate length in constant
time each (via prefix hashes and precomputed powers of the hash base, so no
window needs to be rehashed from scratch), and only pays the cost of a
direct token-by-token comparison when two windows collide on the same hash,
which a large enough hash modulus makes rare.

\begin{algorithm}[h]
\caption{Longest Repeated Substring}
\begin{algorithmic}[1]
\Function{FindLongestRepeatedSubstring}{tokens, minLength}
    \State \textbf{if} tokens too short for two windows of size minLength \textbf{then return} no match
    \State precompute rolling-hash prefix sums over tokens \Comment{$O(n)$, once}
    \Function{ExistsNonOverlappingRepeat}{$L$}
        \State buckets $\gets$ empty map from window hash to start positions
        \For{each start position $i$}
            \State $h \gets$ hash of the length-$L$ window starting at $i$ \Comment{$O(1)$}
            \For{each earlier position $p$ in buckets[$h$] at distance $\geq L$}
                \If{the two windows match token by token}
                    \State \Return $(p, i)$ \Comment{verified match}
                \EndIf
            \EndFor
            \State add $i$ to buckets[$h$]
        \EndFor
        \State \Return not found
    \EndFunction
    \State $lo, hi \gets$ minLength, $\lfloor n/2 \rfloor$; best $\gets$ no match
    \While{$lo \leq hi$}
        \State $mid \gets \lfloor (lo+hi)/2 \rfloor$
        \If{\Call{ExistsNonOverlappingRepeat}{$mid$} found a pair}
            \State best $\gets$ that pair, at length $mid$; $lo \gets mid+1$ \Comment{try longer}
        \Else
            \State $hi \gets mid-1$ \Comment{try shorter}
        \EndIf
    \EndWhile
    \State \Return best
\EndFunction
\end{algorithmic}
\end{algorithm}

The digit-normalized variant runs the same routine on a copy of the
sequence in which every maximal run of digit tokens has been collapsed to a
single wildcard token, then maps the result back to real token positions.

\subsubsection{Interpretation and limitations}
\label{sec:lrs-tolerance}

The model can also vary a letter, a variable name, whitespace, punctuation,
or word choice between iterations, and because LRS only ever looks for
\emph{exact} matches, any such difference can still break a match or shrink
it down to whatever fragment is identical on both sides. Unlike numbers,
there is no fixed, enumerable set of "things that vary" here -- it can be
anything -- so this remains an open limitation of the exact-match approach.

A natural extension would let the matcher tolerate a handful of differing
tokens between two occurrences, not just digits. We chose not to, for three
reasons. First, it makes the algorithm substantially more complex and
computationally heavier: an exact match admits hashing and a clean binary
search over candidate lengths, because two windows are equal exactly when
their hashes are equal, so every comparison is an $O(1)$ hash lookup. An
approximate match has no such shortcut: a single differing token changes a
rolling hash completely, regardless of how many other tokens still agree,
so hash equality carries no information about how close two windows are,
and every candidate pair would instead need an explicit,
token-by-token comparison. Second, unlike the decision thresholds we
calibrate elsewhere in this work (a repeat-count cutoff for LRS itself, a
repetition-score cutoff for TTR, an entropy-drop cutoff for Entropy,
\S\ref{sec:empirical-performance}), a token tolerance would not be a
threshold applied on top of an already-defined score: it would change what
counts as a match in the first place, before any score exists to threshold,
so it is not the same kind of parameter as the ones we do calibrate. Third,
and for the same reason, it has no natural, unit-independent scale to
calibrate against: a tolerance of one differing token is enormous relative
to a four-token repeating unit and negligible relative to a four-hundred-token
one, so the "right" tolerance would have to vary with the length of
whatever unit is being matched, not just with the model or domain, in a way
a single scalar threshold does not. We therefore kept LRS an exact-match
algorithm, with digit normalization as the only, narrowly-motivated
exception, reserving calibration for the classification threshold applied
to LRS's output rather than for the matching rule itself.

How often LRS's tuned threshold fires on naturally-terminated versus
token-capped rollouts, alongside the same check for TTR and Entropy, is
reported in \S\ref{sec:empirical-performance}.

\subsection{LLM Judge}
\label{sec:llm-judge-labeling}

\subsubsection{Setup}

We use \textbf{Claude Sonnet 5} as the judge, called either directly through
the Anthropic API with structured JSON output, or through the Claude Code
CLI (drawing on subscription usage rather than metered billing). Both paths
use the same system prompt, user prompt, and output schema, so results are
directly comparable regardless of which one produced them.

\subsubsection{System prompt and output}

The judge sees only the prompt and the completion text, never the
heuristic scores above, so its verdict is not biased by them:

\begin{quote}
\small\ttfamily\raggedright
You are judging whether a language model's completion degenerated into repetitive, low-information text (looping on a phrase, sentence, or pattern without making progress), as opposed to repetition that is a normal, expected part of the task. You will see both the prompt and the completion -- read the prompt first, since it changes what counts as degeneration.

Two things that look repetitive but are NOT degeneration:

1. Genuine incremental work. Brute-force search, case enumeration, or a step-by-step derivation often produces a fixed template with a changing value each time (e.g. "N = 41: ..." then "N = 42: ..." then "N = 43: ..."). This is fine as long as each occurrence is doing new work -- checking a new candidate, advancing to the next step, covering a new case. A useful test: if you deleted every occurrence but the first, would the completion lose information needed to solve the task? If yes, it's legitimate enumeration, not degeneration. Only flag it if the "progress" is illusory -- e.g. it starts re-checking values it already covered, drifts into a template that stops advancing, or cycles without ever converging on an answer.

2. Instructed repetition. If the prompt itself asks for repeated output (e.g. "repeat X ten times", "write a line for each of the following ..."), producing that repetition is correct instruction-following, not degeneration -- unless the completion overruns what was asked (keeps going well past the requested count) or the repetition itself becomes corrupted or loops in a way the prompt didn't ask for.

Completions are all hard-capped at the same fixed max-token budget. A completion that ends abruptly mid-sentence, mid-word, or mid-token is expected simply because it hit that cap -- this is not, by itself, evidence of degeneration one way or the other. Judge based on the quality and content of the text alone, not on whether or how it was cut off.
\end{quote}

The judge returns, as structured JSON, a \texttt{reasoning} field (written
\emph{before} the verdict, so it acts as genuine scratch space rather than a
post-hoc justification), the boolean verdict \texttt{is\_degenerating}, an
\texttt{onset\_quote}, a short verbatim quote marking where the pattern
\emph{first} begins, if degenerating, and a \texttt{confidence} level.
Before writing \texttt{onset\_quote} into its final answer, the judge must first call a \texttt{verify\_onset\_quote}
tool with its candidate quote. The tool performs a plain verbatim substring search against the real completion text
plus a minimum-length check, and reports back whether the candidate was
found, giving the judge one chance to fix a bad quote before finalizing its
answer.

\subsubsection{Why we trust it}

The LLM judge is the only metric we have that reliably tells apart real
degeneration from text that merely \emph{looks} repetitive by construction,
the clearest case being brute-force problem solving, which is correctly,
legitimately repetitive (the same template, a new value checked each time),
precisely because genuine exhaustive enumeration can need far more tokens
than our fixed generation budget allows and so is truncated identically to
a real loop, with no way to tell the two apart from the stop reason alone.
On the small set of judge-confirmed-negative, token-capped rollouts (the
population this case corresponds to), TTR's tuned threshold fires on the
large majority of them and LRS's fires on a build-dependent fraction
(\S\ref{sec:hard-negatives}). This is exactly what point 1 of the system prompt
above is written to handle: because the judge understands the
\emph{semantics} of the reasoning, not just the surface repetition of
tokens, it can tell a loop that keeps re-deriving the same thing apart from
one that is actually making progress case by case.

\subsubsection{Cost and feasibility}

We only send the judge rollouts that hit the $4096$-token generation cap
without reaching a natural stop, the population where degeneration
overwhelmingly concentrates: among the rollouts we send it, the judge
independently confirms degeneration in 93.0--99.5\% of cases across the
three dataset builds (97.1\% pooled), i.e.\ \texttt{stop\_reason ==
"length"} alone is already a strong, if imperfect, proxy for real
degeneration. This is a small slice of each dataset: 890/36{,}300
($\sim$2.5\%) rollouts for Apertus-8B-Instruct, 1348/36{,}300 ($\sim$3.7\%) for Apertus 1.5
capfilter-linear, and 1193/36{,}300 ($\sim$3.3\%) for Apertus 1.5 SFT-256k.
Keeping the judged population this small keeps the judging cost
affordable, and, just as important, keeps the resulting set small
enough to manually read through and spot-check individual verdicts, which
would not be realistic if every rollout were judged.

\subsubsection{Reliability}
\label{sec:judge-reliability}

\emph{\texttt{onset\_quote} verbatim rate.} Of the \texttt{onset\_quote}
values the judge returns, 99.6--100\% are found verbatim in the completion
text they were judging, checked independently of the judge's own
\texttt{verify\_onset\_quote} tool call.

One reliability gap is worth noting explicitly: \texttt{onset\_quote} can
contain a generic placeholder that coincidentally matches unrelated text
elsewhere in the completion, observed once as the literal word
\texttt{"test"}, which matched an unrelated "ratio test" phrase in a
convergence problem. The pipeline guards against this with two checks: the
\texttt{verify\_onset\_quote} tool lets the judge itself catch and fix an
under-four-word candidate before finalizing its answer, and, since that
in-context check only ever sees whatever candidate the model chooses to run
it on, a binding server-side check independently re-verifies the
\emph{final} \texttt{onset\_quote} field against the real completion text on
the same two criteria (length and verbatim presence) before a degenerating
verdict is accepted. A row that fails the server-side check is re-judged
automatically on the next run rather than accepted with a bad
\texttt{onset\_quote}.

A second, unrelated failure mode is a hard error rather than a bad verdict:
a small fraction of rollouts are refused outright by the Claude Code CLI
path, subject to its own usage-policy screening (separate from the raw
Anthropic API). Across all three dataset builds this affects 211/5674
($\sim$3.7\%) of Claude-Code-CLI judge calls, overwhelmingly concentrated in
one domain, 91.5\% of refusals come from \texttt{if\_sft\_data\_verified},
whose prompts often carry unusual formatting constraints (e.g. "respond
entirely in capital letters", "include at least $N$ placeholder tokens")
that, combined with degenerate output, can resemble spam to a usage-policy
classifier. These rows are manually judged rather than retried.

\subsection{Empirical Performance}
\label{sec:empirical-performance}

Entropy, TTR, and LRS are evaluated below against one shared population and
protocol. \textbf{Positives} are rollouts that hit the $4096$-token
generation cap and receive an \texttt{is\_degenerating=True} verdict from
the LLM judge (\S\ref{sec:llm-judge-labeling}). \textbf{Negatives} are every
rollout that reaches a natural end-of-sequence, by definition, with no
judge call needed -- \S\ref{sec:flag-rate-by-termination} shows that all
three metrics fire far more often on token-capped rollouts than on
naturally-terminated ones, so treating end-of-sequence as a safe negative
label costs essentially nothing while removing the need to judge tens of
thousands of additional rollouts per dataset build. The small number of
judge-confirmed-negative rollouts among the token-capped population are
also kept, as informative hard negatives. For every metric and dataset
build, a decision threshold (a repeat-count cutoff for LRS, a score cutoff
for TTR and Entropy) is tuned once by maximizing balanced accuracy on one
deterministic half of the prompts (the calibration set), and every number
reported below comes from the disjoint other half (the test set).

\subsubsection{Aggregate results}

Every number in this section is at the rollout level, not the token level:
each rollout is one confusion-matrix instance, classified positive or
negative by its own metric's already-defined rollout-level decision rule
($r_{\max}>\tau$ for TTR, the analogous entropy-drop-score cutoff for
Entropy, a repeat-count cutoff for LRS), evaluated at the threshold given
in Table~\ref{tab:unified-aggregate}. Table~\ref{tab:unified-aggregate}
reports, for every metric and dataset
build, the domain-balanced macro-average of the completion-level confusion
matrix: each domain's score is computed separately and then averaged with
equal weight, rather than pooling all rollouts into one confusion matrix
first. Pooling would let whichever domain happens to contribute the most
rollouts dominate the result; macro-averaging instead treats every domain
as equally informative regardless of its size. Two of the seven domains
(\texttt{aime\_2025},
\texttt{codeforces}) contribute fewer than $10$ positive test rollouts for
every metric on at least one build, and a third (\texttt{medical\_o1})
contributes none at all on two of the three builds -- its single positive
rollout, out of $6000$, happens to fall on the calibration side of the
split for those builds by chance. These domains are always reported in the
per-domain table below, but excluded from the macro-average here, since a
metric computed on $0$--$8$ positive examples carries no statistical
weight.

\begin{table}[h]
\centering
\resizebox{\textwidth}{!}{%
\begin{tabular}{llccccccc}
\toprule
Metric & Dataset build & Threshold & Precision & Recall & Specificity & Bal. Acc. & MCC & Accuracy \\
\midrule
Entropy & Apertus-8B-Instruct             & $0.4835$    & 0.061 & 0.731 & 0.529 & 0.630 & 0.103 & 0.535 \\
Entropy & Apertus 1.5 -- capfilter-linear & $0.4505$    & 0.109 & 0.623 & 0.717 & 0.670 & 0.162 & 0.718 \\
Entropy & Apertus 1.5 -- SFT-256k         & $0.4837$    & 0.085 & 0.640 & 0.737 & 0.689 & 0.156 & 0.736 \\
TTR     & Apertus-8B-Instruct             & $0.55$      & 0.157 & 0.717 & 0.840 & 0.778 & 0.277 & 0.835 \\
TTR     & Apertus 1.5 -- capfilter-linear & $0.65$      & 0.330 & 0.897 & 0.921 & 0.909 & 0.494 & 0.920 \\
TTR     & Apertus 1.5 -- SFT-256k         & $0.60$      & 0.300 & 0.787 & 0.871 & 0.829 & 0.376 & 0.870 \\
LRS     & Apertus-8B-Instruct             & $\geq 3$    & 0.388 & 0.834 & 0.957 & 0.895 & 0.547 & 0.953 \\
LRS     & Apertus 1.5 -- capfilter-linear & $\geq 4$    & 0.595 & 0.915 & 0.978 & 0.946 & 0.718 & 0.976 \\
LRS     & Apertus 1.5 -- SFT-256k         & $\geq 3$    & 0.377 & 0.812 & 0.955 & 0.884 & 0.520 & 0.951 \\
\bottomrule
\end{tabular}%
}
\caption{Domain-balanced macro-average of completion-level detection
performance, test split only, for all three structural metrics against the
LLM judge. Threshold is the entropy-drop-score cutoff (Entropy),
repetition-score cutoff (TTR), or repeat-count cutoff (LRS) tuned per
metric and dataset build on a disjoint calibration split
(\S\ref{sec:empirical-performance}).}
\label{tab:unified-aggregate}
\end{table}

Precision is low for every metric under this realistic population --
$0.061$--$0.109$ for Entropy, $0.157$--$0.330$ for TTR, $0.377$--$0.595$ for
LRS -- because natural end-of-sequence rollouts vastly outnumber
token-capped ones in a real dataset, so even a handful of false positives
per thousand true negatives dominates the denominator. Ranked by balanced
accuracy and MCC, LRS is the strongest standalone detector on every build,
TTR is a clear second, and Entropy is the weakest, most exploratory signal:
its specificity ($0.529$ for Apertus-8B-Instruct) shows that a persistent
entropy drop is common in ordinary, non-degenerating text, consistent with
its exploratory role in \S\ref{sec:entropy-labeling}.

\subsubsection{Per-domain results}

Table~\ref{tab:unified-per-domain} breaks the same comparison down by
domain, for Apertus-8B-Instruct. $n$ is the number of test-split positive
rollouts in that domain -- identical across all three metrics, since it
depends only on the judge's verdict, never on any metric's own threshold.
Domains marked $\dagger$ fall under the $10$-positive-rollout line used for
the macro-average above and should be read as anecdotal rather than a
generalizable estimate; \texttt{medical\_o1} has none at all on this build.

\begin{table}[h]
\centering
\resizebox{\textwidth}{!}{%
\begin{tabular}{lccccccc}
\toprule
& & \multicolumn{2}{c}{Entropy} & \multicolumn{2}{c}{TTR} & \multicolumn{2}{c}{LRS} \\
\cmidrule(lr){3-4}\cmidrule(lr){5-6}\cmidrule(lr){7-8}
Domain & $n$ & Precision & Bal. Acc. & Precision & Bal. Acc. & Precision & Bal. Acc. \\
\midrule
aime\_2025$^\dagger$    & 8   & 0.055 & 0.536 & 0.148 & 0.849 & 0.241 & 0.865 \\
codeforces$^\dagger$    & 7   & 0.004 & 0.635 & 0.031 & 0.962 & 0.053 & 0.978 \\
deepmath\_103k          & 60  & 0.036 & 0.624 & 0.066 & 0.774 & 0.299 & 0.935 \\
if\_sft\_data\_verified & 135 & 0.116 & 0.678 & 0.241 & 0.776 & 0.459 & 0.910 \\
llama\_nemotron         & 65  & 0.029 & 0.622 & 0.199 & 0.773 & 0.373 & 0.804 \\
medical\_o1$^\dagger$   & 0   & --    & --    & --    & --    & --    & -- \\
numinamath\_1\_5        & 116 & 0.063 & 0.597 & 0.122 & 0.791 & 0.422 & 0.932 \\
\bottomrule
\end{tabular}%
}
\caption{Per-domain precision and balanced accuracy, Apertus-8B-Instruct,
test split. $^\dagger$: fewer than $10$ positive test rollouts, excluded
from Table~\ref{tab:unified-aggregate}'s macro-average.}
\label{tab:unified-per-domain}
\end{table}

Precision and balanced accuracy are shown, rather than the full six-metric
set of Table~\ref{tab:unified-aggregate}, because together they isolate the
two effects of interest here: precision shows how much each domain's
population size drags down the raw hit rate, while balanced accuracy shows
that the same ranking holds regardless of population size. The domain breakdown shows the aggregate ranking is not driven by one
unusually easy or unusually hard domain: LRS's balanced accuracy exceeds
TTR's, which in turn exceeds Entropy's, in every one of the seven domains,
not just in the pooled or macro-averaged view.

\subsubsection{Threshold-free comparison}

Figure~\ref{fig:roc-comparison} sweeps every possible decision threshold
for each metric and reports the resulting ROC curve and its area under the
curve (AUC), as a sanity check that the single calibrated operating point
above was not a lucky pick. The ranking is consistent with the
operating-point results: LRS has the highest AUC on every build ($0.943$,
$0.980$, $0.922$), TTR is second ($0.852$, $0.951$, $0.894$), and Entropy is
last and closest to the diagonal ($0.695$, $0.732$, $0.776$).

\begin{figure}[h]
\centering
\includegraphics[width=\textwidth]{figures/roc_comparison.pdf}
\caption{Threshold-free ROC curves for Entropy, TTR, and LRS against the
LLM judge, test split only, one panel per dataset build.}
\label{fig:roc-comparison}
\end{figure}

\subsubsection{Onset agreement}

For rollouts where a metric predicts degeneration and the judge's
\texttt{onset\_quote} is located verbatim in the completion,
Table~\ref{tab:unified-onset} reports the signed distance between the
judge's onset and the metric's own onset (judge onset minus metric onset:
positive means the metric fires earlier than the judge). All three
onset-producing rules are now compared on the same underlying positive
population for the first time, since they all draw from the same
judge-confirmed positives (\S\ref{sec:empirical-performance}) rather than
three separately constructed ones.

\begin{table}[h]
\centering
\resizebox{\textwidth}{!}{%
\begin{tabular}{llcccc}
\toprule
Metric & Dataset build & $n$ & Median signed & Median abs. & IQR (signed) \\
\midrule
Entropy     & Apertus-8B-Instruct             & 574  & 283.0 & 339.5 & [$-11.0$, $703.8$] \\
Entropy     & Apertus 1.5 -- capfilter-linear & 881  & 141.0 & 168.0 & [$-24.0$, $559.0$] \\
Entropy     & Apertus 1.5 -- SFT-256k         & 689  & 133.0 & 253.0 & [$-41.0$, $725.0$] \\
TTR (run-4) & Apertus-8B-Instruct             & 622  & 274.0 & 311.0 & [$102.0$, $680.5$] \\
TTR (run-4) & Apertus 1.5 -- capfilter-linear & 1139 & 160.0 & 165.0 & [$91.0$, $339.5$] \\
TTR (run-4) & Apertus 1.5 -- SFT-256k         & 846  & 236.5 & 267.0 & [$115.3$, $771.8$] \\
LRS         & Apertus-8B-Instruct             & 815  & 1.0   & 106.0 & [$-189.0$, $78.5$] \\
LRS         & Apertus 1.5 -- capfilter-linear & 1267 & 1.0   & 29.0  & [$-24.0$, $32.0$] \\
LRS         & Apertus 1.5 -- SFT-256k         & 1031 & 1.0   & 70.0  & [$-98.5$, $57.0$] \\
\bottomrule
\end{tabular}%
}
\caption{Signed onset distance (judge onset minus metric onset, in tokens)
for rollouts where both are defined, test split only.}
\label{tab:unified-onset}
\end{table}

LRS's onset agrees with the judge almost exactly on median (within a single
token on every build), reflecting that both are, in different ways,
locating the same first exact (or digit-normalized) repeat. TTR and Entropy
instead fire well before the judge's chosen onset on median (roughly
$130$--$280$ tokens early), consistent with each detecting an earlier, more
gradual signal -- a sustained rise in local repetition or a sustained drop
in entropy -- rather than the specific point a human reader would call
unambiguous. All three distributions remain wide (interquartile ranges
spanning several hundred tokens in most cases), so none of the three
metrics should be treated as pinpointing onset precisely on any individual
rollout; they are more useful in aggregate than case by case. Median
absolute distance is reported alongside the signed one because a near-zero
median signed distance does not by itself imply pinpoint accuracy: LRS's
median signed distance is within a single token of zero on every build, yet
its median absolute distance is still $29$--$106$ tokens, showing that its
error is unbiased in direction but not negligible in typical size.

\subsubsection{Flag rate by termination reason}
\label{sec:flag-rate-by-termination}

As a sanity check on treating natural end-of-sequence as a safe negative
label, Table~\ref{tab:flag-rate-termination} reports, for every metric and
dataset build, the fraction of rollouts each metric's already-tuned
threshold flags as positive, split by whether the rollout reached
end-of-sequence or hit the token cap -- computed over every rollout in the
dataset, not just the judge-touched subset used elsewhere in this section,
since no judge call is needed to know a rollout's own termination reason.
For all three metrics the gap between the two columns is large, confirming
that a metric firing is a rare event on naturally-terminated text and a
common one on token-capped text -- exactly the assumption the population
construction above depends on. Entropy's EOS flag rate ($20.5$--$38.6\%$) is
an order of magnitude higher than TTR's ($7.0$--$13.3\%$) or LRS's
($2.1$--$4.8\%$), consistent with it being the noisiest of the three
signals throughout this section.

\begin{table}[h]
\centering
\resizebox{\textwidth}{!}{%
\begin{tabular}{llcc}
\toprule
Metric & Dataset build & EOS flag rate & Token-capped flag rate \\
\midrule
Entropy & Apertus-8B-Instruct             & 38.6\% & 69.7\% \\
Entropy & Apertus 1.5 -- capfilter-linear & 28.6\% & 69.4\% \\
Entropy & Apertus 1.5 -- SFT-256k         & 20.5\% & 67.6\% \\
TTR     & Apertus-8B-Instruct             & 12.1\% & 77.3\% \\
TTR     & Apertus 1.5 -- capfilter-linear & 7.0\%  & 89.6\% \\
TTR     & Apertus 1.5 -- SFT-256k         & 13.3\% & 82.4\% \\
LRS     & Apertus-8B-Instruct             & 3.9\%  & 88.9\% \\
LRS     & Apertus 1.5 -- capfilter-linear & 2.1\%  & 92.4\% \\
LRS     & Apertus 1.5 -- SFT-256k         & 4.8\%  & 75.4\% \\
\bottomrule
\end{tabular}%
}
\caption{Fraction of rollouts flagged positive by each metric's own tuned
threshold (\S\ref{sec:empirical-performance}), split by termination reason,
over every rollout in the dataset.}
\label{tab:flag-rate-termination}
\end{table}

\subsubsection{Flag rate on hard negatives}
\label{sec:hard-negatives}

Restricting further to the judge-confirmed-negative rollouts among the
token-capped population -- the small set of "hard negatives" kept in the
unified population (\S\ref{sec:empirical-performance}), and the population
\S\ref{sec:llm-judge-labeling}'s brute-force-solving argument refers to --
Table~\ref{tab:hard-negatives} reports how often each metric's tuned
threshold still fires on them. This population is small (6--78 rollouts
per build, since a rollout must both hit the token cap and be confirmed
non-degenerate by the judge), so the two smallest builds should be read
with that caveat in mind. TTR fires on the large majority of hard
negatives on every build ($79.5$--$100\%$), supporting the idea that
legitimate, repetitive-looking text such as brute-force enumeration
reliably trips a local repetition metric. LRS's tuned threshold fires far
less consistently ($7.7$--$75.0\%$), and Entropy's shows no clear pattern
($12.5$--$87.2\%$) -- both are noisier indicators of this specific failure
mode than TTR.

\begin{table}[h]
\centering
\resizebox{\textwidth}{!}{%
\begin{tabular}{lcccc}
\toprule
Dataset build & $n$ & LRS flag rate & TTR flag rate & Entropy flag rate \\
\midrule
Apertus-8B-Instruct             & 8  & 75.0\% & 100.0\% & 12.5\% \\
Apertus 1.5 -- capfilter-linear & 6  & 33.3\% & 100.0\% & 83.3\% \\
Apertus 1.5 -- SFT-256k         & 78 & 7.7\%  & 79.5\%  & 87.2\% \\
\bottomrule
\end{tabular}%
}
\caption{Fraction of judge-confirmed-negative, token-capped rollouts
("hard negatives") each metric's own tuned threshold still flags as
positive.}
\label{tab:hard-negatives}
\end{table}

\subsubsection{Evaluation metric definitions}

For a completion-level confusion matrix with true positives (TP), false
positives (FP), true negatives (TN), and false negatives (FN):
\[
\text{Precision}=\frac{TP}{TP+FP},
\qquad
\text{Recall}=\frac{TP}{TP+FN},
\qquad
\text{Specificity}=\frac{TN}{TN+FP},
\]
\[
\resizebox{0.92\linewidth}{!}{$\text{Balanced accuracy}=\dfrac{\text{Recall}+\text{Specificity}}{2},
\qquad
\text{Accuracy}=\dfrac{TP+TN}{TP+FP+TN+FN}.$}
\]
\[
\resizebox{0.92\linewidth}{!}{$\text{MCC}=\dfrac{TP\cdot TN-FP\cdot FN}{\sqrt{(TP+FP)(TP+FN)(TN+FP)(TN+FN)}}.$}
\]
The Matthews correlation coefficient (MCC) folds all four confusion-matrix
cell counts into a single number in $[-1,1]$, with $0$ corresponding to a
rule no better than chance and $1$ to perfect prediction; unlike accuracy,
it stays informative even when positives are rare, which is the case
throughout this section.

\end{document}
